\section{Consciousness Explained}

\textbf{Daniel Dennett}, a professor of philosophy at my alma mater, is one of the ``New Atheists''. Several years ago, a small group of friends use to meet to discuss esoteric texts. One fellow began to ``lose his religion'' and announced that the next text would be \emph{Consciousness Explained} followed by \emph{Darwin's Dangerous Idea}, both by Dennett.

Consciousness explained was ultimately anti-climactic since the ending was more like consciousness explained away. On the principle that the errors of an intelligent man are more interesting than the truths of a mediocrity, there is much value in Dennett's books. First of all, they provide the latest researches into neuroscience and biology, integrated into his peculiar worldview.

\paragraph{Naturalism}


Naturalism is the view that the world process can be fully described solely in terms of natural, i.e., physical, causes. In other words, the cosmos is epistemologically closed. The cosmos can be explained from internal evidence alone, requiring no intervention by supernatural or preternatural causes. In logical terms, that means that the system is complete. Unfortunately, in actuality it is not at all complete, since there are things and events that have no plausible explanation, hence the search for the ``Theory of Everything.'' Nevertheless, naturalists take it on faith that these explanations will be found.

A logical system, moreover, needs to be consistent so that no self-contradictions can arise. We don't know that, and there are reasons to believe that it is not consistent.

Naturalism must necessarily reject two seemingly self-evident truths: first, the idea of God, and second, more startingly, the inefficacy of consciousness.

From a metaphysical perspective, naturalism recognizes matter, space, and time as the sole conditions of manifestation. It leaves out essence and psyche.

\paragraph{The Notion of God}


From that perspective, the notion of God serves no purpose. In other words, the theories of physics, the laws of mathematics, chemical reactions, and biological development can be studied without reference to any God. I've heard enough debates with physicists who simply look puzzled when God is invoked. By Occam's razor, ``God'' is a superfluous hypothesis, not to mention that the description of this God is quite muddled, and varies quite a bit from believer to believer.

The point is that the idea of God is incompatible with naturalism, and is always embedded within a more complex metaphysical framework.

\paragraph{Consciousness}


Consciousness is required by no physical equation nor any scientific theory. At best, it is a useless epiphenomenon, having no power to affect physical events. Dennett has to go through pains to demonstrate the lack of free will. For example, he refers to neurological tests the predict the choices of a subject even before he becomes consciously aware of having made that choice.

The idea of Qualia, he claims, is too vague to be of value. The subjective sensation of color, for example, changes nothing; the frequency of electromagnetic waves is all that matters. He provides countless examples.

There is no point being concerned about conscious experience, although everyone is. All your pains, anxiety, depression, fears, and so on, do not matter for the world process. Neither do your joys, pleasures, and so on. They are merely epiphenomena of biological processes.

That view is eerily reminiscent of those Gnostic sects that believed we were created by an evil demiurge. The Demiurge trapped our consciousnesses in a world of matter, with no ability to escape. We are doomed to suffer in matter, yet unable to do anything about it.

\paragraph{Moralism}


That Gnostic view is quite terrifying, yet is has no such effect on Dennett who seems quite sanguine about the value of his philosophy. Moreover, he inconsistently hangs onto the idea that there is a moral way to act in the world. Not surprisingly, his views on that, seem to match about 90\% of the academic community. So Dennett moralizes, trying to spread atheism far and wide. For what purpose, we don't know; at least, it is not clear to me. A belief should also be an epiphenomenon with no causal efficacy in the world process.

\textbf{George Santayana} is a much deeper thinker and more consistent that Dennett. For him, there is no point to promote his philosophy. Those who can understand, will understand; otherwise, it is pointless to explain it to the ignorant.

\paragraph{A Day in the Life}


Let's pretend that naturalism is true. We can analyze a typical day in a small Potemkin like village, in which the human beings can be hooked up to neuronal devices and other medical devices such as oxygen intake and blood pressure.

We can verify that energy and momentum are conserved in the village. We can measure the caloric intake and expenditures to verify it. We know how their physiology functions. We watch them move about the village and can measure the expenditure of energy.

So far no laws of physics have been violated. The people of the village can be studied as if they were billiard balls on a pool table.

Their human interactions can be explained by evolutionary psychology. They bond with their closest genetic relatives and treat strangers with more caution. Since humans are social, we observe how social hierarchies are established, and so on.

Nevertheless, there are some inconsistencies that demand explanation. By Newton's law, motion should remain in a straight line unless acted upon by an outside force. Yet humans never walk straight for very long, making sudden uncaused changes of direction.

The force, then, must be interior, not exterior. We can see through our instruments that each change of direction is preceded by some neuronal activity. That is the sought for cause, even though it explains nothing.

If the human being turns to enter a grocery store, we can understand that it is seeking a new energy source. But if it walks into a comic book store, the event seems to have not biologically necessary cause.

TO BE CONTINUED


\flrightit{Posted on 2021-07-08 by Cologero}

